\documentclass[10pt,a4paper]{amsart}
\usepackage[margin=19mm]{geometry}
\usepackage{amsmath,amssymb,mathtools}
\usepackage[scr=rsfs,cal=euler]{mathalfa}
\usepackage{xfrac}
\usepackage[unicode,hidelinks,bookmarks=false]{hyperref}
\newcommand{\ie}{i.e.\ }
\newcommand{\cf}{cf.\ }
\newcommand{\resp}{respectively\ }
\newcommand{\Subsection}{\subsection}
\providecommand{\nobreakdash}{\nobreak\hbox{-}}
\setlength{\emergencystretch}{2em}
\multlinegap0pt
\theoremstyle{plain}
\newtheorem{theorem}{Theorem}
\newtheorem{proposition}[theorem]{Proposition}
\newtheorem{lemma}[theorem]{Lemma}
\newtheorem{corollary}[theorem]{Corollary}
\theoremstyle{definition}
\newtheorem{definition}[theorem]{Definition}
\theoremstyle{remark}
\newtheorem{remark}[theorem]{Remark}
\title[Dynamic Interval Scheduling with Random Start and End Times]{Dynamic Interval Scheduling with Random Start and End Times: NP-hardness of the CDSRSE model (printed as Theorem 3.2)}
\author{Rui Gong, Alejandro Toriello}
\date{Source: arXiv:2602.07217v3 (CC BY 4.0) (CC BY 4.0)}
\begin{document}
\maketitle

\noindent\emph{Source and licence.} Dynamic Interval Scheduling with Random Start and End Times. Version arXiv:2602.07217v3 (CC BY 4.0), distributed under the Creative Commons Attribution 4.0 International licence, http://creativecommons.org/licenses/by/4.0/, as recorded in the arXiv licence metadata for this exact version and shown on the saved abstract page https://arxiv.org/abs/2602.07217v3. Full licence notice: licenses/arxiv-2602.07217-CC-BY-4.0.txt.\par\smallskip

\noindent\textit{Editorial scope.} Counted unit: the NP-hardness of the CDSRSE scheduling model, printed in the article as Theorem 3.2 (source0.tex lines 418--420, one \texttt{theorem} environment), delivered together with the article's complete original proof of it (source0.tex lines 421--443, one \texttt{proof} environment of 23 lines immediately adjacent to the statement, with no gap and no deferral). The proof is the article's own reduction: it builds the instance with one long task and n singleton tasks, verifies that under the model the second task is always deleted because its slot is occupied, and then reduces the decision problem to the dynamic maximum stable set problem on a star graph, which is known to be NP-hard. No mathematics is written, repaired or re-derived: statement and proof are the article's own text line for line, in the article's own notation, and no line of either is omitted, substituted or re-flowed.

Required context retained uncounted: the article's statement that the DSRSE model is NP-hard, with its label (source0.tex lines 248--250), which the delivered proof invokes as its starting point. Nothing else from the article is needed.

References. Two citation commands occur in the retained lines: the citation of the dynamic node packing problem and the citation of the star-graph NP-hardness. The package therefore carries exactly those two bibliography entries, transcribed from the article's own BibTeX database, printed with the article's own labels 8 and 1.

Editorial formatting only, in addition: an AMS \texttt{amsart} document that restates the article's own theorem-environment declarations with the same shared counter the article declares. No macro definition is needed by any retained line, so none is restated. The article's own class is the standard article class; no publisher class file is shipped or used. Consequently the retained environments print here in the order in which they occur, so the counted theorem prints as Theorem 2 whereas the article prints it as Theorem 3.2 under its per-section numbering, and the retained proposition prints as Proposition 1, which the article prints as Proposition 3.1; the equations of the retained spans are numbered anew in order of appearance and every cross-reference inside the retained text resolves to this wrapper's numbering. That is the only change to the numbering. No figure, no image-inclusion line and no drawing code occur in any retained line, so no artwork is delivered or needed, although the article contains figures elsewhere.

Not retained: the article's abstract, its introduction, its model definitions, its examples, its algorithms and its other results; none of that material is used as a step of the delivered proof, and the delivered proof is complete as the author wrote it. The package ships the article's concatenated source verbatim as \texttt{sources/194159/source0.tex}, and every delivered excerpt is an exact line range of that file. No mathematical statement, hypothesis, estimate or conclusion has been rewritten, repaired or added, and printed slips, if any, are preserved literally.
\begin{proposition}\label{prop:NP-hard}
     The DSRSE model is NP-hard.
\end{proposition}

\begin{theorem}
     The CDSRSE model is also NP-hard.
\end{theorem}

\begin{proof}
    As in the previous proof, consider the instance with $n+1$ tasks and $2n$ slots. 
    Task $0$ occupies slots $1,\ldots,n$ with probability $1$, and task $i$ occupies slot $i$ with probability $p_i$ and $n+i$ with probability $q_i:=1-p_i$. {As in the proof of Proposition \ref{prop:NP-hard}, task $0$ is an interval task and tasks $1,\ldots,n$ are singleton tasks.}
    Following \cite{DNP}, any solution for this problem is determined by when to select task $0$, and what is selected before trying to select task $0$.
    Then the problem can be formulated as
    \begin{equation}
    \label{Conserv-JS}
    \tag{Conserv-JS}
    \max_{S\subseteq N}\left\{\sum_{i\in S}w_i+(1-\prod_{i\in S}q_i)\sum_{j\notin S}w_j+\prod_{i\in S}q_i{w_0}\right\}.
    \end{equation}
    By rearranging, we get 
    \[
    \max_{S\subseteq N}\left\{\sum_{i\in N}w_i+\prod_{i\in S}q_i(w_0-\sum_{j\notin S}w_j)\right\}.
    \]
    {For the positive weights and probabilities used below, we restrict the following logarithmic formulation to subsets with $w_0-\sum_{j\notin S}w_j>0$; choosing $S=N$ gives a positive contribution, so this restriction does not change the optimum.}
    By dropping constant terms and applying the natural logarithm, we get 
    \[\max_{S\subseteq N}\left\{\sum_{i\in S}\ln q_i+\ln(w_0-\sum_{j\notin S}w_j)\right\}.\]
% ATH: You already use a_i above; need different notation
    Given a set of positive numbers $\zeta_i$, $i \in N$, the partitioning problem asks for a subset $S \subseteq N$ such that $\sum_{i\in S} \zeta_i = \sum_{j\notin S} \zeta_j$; without loss of generality, we assume that $\sum_{i\in N} \zeta_i = 2$.
    By setting parameters of the problem above as $q_i = e^{- \zeta_i}$, $p_i = 1-e^{-\zeta_i}$, $w_i =\zeta_i$,and $w_0 = \sum_{i\in N} \zeta_i$, we formulate
    \[\max_{{\emptyset\ne S\subseteq N}}\left\{-\sum_{i\in S} \zeta_i+\ln(\sum_{i \in S} \zeta_i)\right\}.\]
Proceeding similarly to~\cite{Star-Graph-NP-hard}, there is a set $S \subseteq N$ with $\sum_{i\in S} \zeta_i=1$ if and only if the optimum of this problem is $-1$.
\end{proof}

\begin{thebibliography}{99}
\bibitem[8]{DNP} Christopher Muir and Alejandro Toriello, Dynamic node packing, Mathematical Programming 196 (2022), 875--906.
\bibitem[1]{Star-Graph-NP-hard} Shabbir Ahmed and Alper Atamt\"{u}rk, Maximizing a class of submodular utility functions, Mathematical Programming 128 (2011), 149--169.
\end{thebibliography}
\end{document}
